\section{Introduction}
\subsection{Work Streams}

% ── Stream A (Pipeline Evaluation) ── DONE ────────────────────────────
% [DONE] Stream A: Pipeline Evaluation
% - Evaluated existing pipeline with multiple hypothesis sets
% - Results shared with @somesh and @harini for review
% - Standalone hypothesis_generator.py (Mode 8) implemented and validated:
%     5 operators x 1 model (GPT-5.5), parallel generation confirmed working
%     Stages 0-5 (dataset context, retrieval, findings, generation,
%     operationalization check, novelty gate) all exercised end-to-end
%     Checkpoint/resume infrastructure and caching (retrieval, findings,
%     embeddings) confirmed functional
% ──────────────────────────────────────────────────────────────────────

% ── Stream B (Literature Loop / AI Co-scientist) ── ACTIVE ───────────
\todo{Stream B: Add literature loop from AI Co-scientist paper; work on iterative feedback for idea refining}
% ──────────────────────────────────────────────────────────────────────

% ── Stream C (Multi-LLM / Model Fusion) ── ACTIVE ────────────────────
\todo{Stream C: Multi-LLM / Model Fusion for hypothesis generation}
% Code context (hypothesis\_generator.py):
%   - Current implementation: Mode 8 = 5 OPERATORS x 1 model (GPT-5.5), 1 persona
%     Each operator (Extension, Combination, Boundary Condition, Mediator, Moderator)
%     runs in parallel via ThreadPoolExecutor(max\_workers=5)
%   - OPERATORS dict and HypothesisGeneratorAgent._call() already support swapping models
%   - NOT YET: true multi-LLM fusion (different model backends per operator call)
%   - NOT YET: model fusion / mixture-of-experts aggregation across outputs
%   - Ref: https://openrouter.ai/blog/announcements/fusion-beats-frontier/
%   - Variants to implement: (a) 5 LLMs x 1 persona, (b) 1 LLM x 5 personas,
%     (c) 5 LLMs x 5 personas — placeholders for ablation table
% ──────────────────────────────────────────────────────────────────────

% ── Stream D (Multi-Agent Social Simulation) ── ACTIVE ───────────────
\todo{Stream D: Multi-agent setup for social simulation, idea, and hypothesis generation}
% Code context (hypothesis\_generator.py):
%   - Current: single-agent generator loop, no inter-agent communication
%   - HypothesisGeneratorAgent is self-contained; no Analyst agent is wired in
%     (stripped from agents\_patched.py per docstring comment)
%   - NOT YET: multi-agent social simulation (debater/critic/proposer roles)
%   - NOT YET: iterative agent-to-agent feedback rounds
%   - Ref: https://behavior-in-the-wild.github.io/social-agents.html
%   - Starting point: wire Analyst agent back in from experigen/algorithm/agents.py
% ──────────────────────────────────────────────────────────────────────

% ── Stream E (Training LLM / Agents) ── ACTIVE ───────────────────────
\todo{Stream E: Training / fine-tuning LLMs or agents for hypothesis generation or operationalization tasks}
% Code context (hypothesis\_generator.py):
%   - Pipeline produces labeled output: hypothesis + op\_pass + novelty\_label per run
%     (saved to results/hypothesis\_candidates.jsonl)
%   - This output is the natural supervised signal for fine-tuning a generator
%   - NOT YET: fine-tuning loop, RLHF/DPO setup, or distillation from GPT-5.5
%   - NOT YET: training operationalization checker (Stage 4) as a smaller model
%   - novelty\_scorer\_v2.py (KNN + embedding + LLM) could serve as reward model
% ──────────────────────────────────────────────────────────────────────

% ── Stream F (Multi-Agent Collaboration Papers) ── DONE ───────────────
% [DONE] Stream F: Multi-Agent Collaboration Papers (@harini's 3 papers)
% - Read all 3 papers provided by @harini on multi-agent collaboration
% - Implemented one approach on the current CMV problem
%   (5 parallel operator calls in generate\_candidates() serve as the
%    multi-agent collaboration scaffold; ThreadPoolExecutor drives them)
% ──────────────────────────────────────────────────────────────────────
    
\subsection{Lit Review for Literature Agents}    

\subsection{Lit Review for AI Co-scientists}

\subsection{Generating Novel and Plausible Ideas}
\subsubsection{Novel}
Already solved to some extent using current algo

\subsubsection{Plausible}
\paragraph{Literature}
1. Authority is persuasive [Persuasion experiment where researchers used electric shocks]
2. Using links or citations is persuasive [from debates]

3. Using links or citations with authorative signals 
[Journals, Nature Papers, and top conferences, .gov websites >>> youtube, wikipedia, or reddit]
% \paragraph{Data}

\paragraph{Agents}
https://behavior-in-the-wild.github.io/social-agents.html